\SourceLeaf{300}
\begin{flushright}
Göttingen, 23 July 1808.
\end{flushright}

\par\indent At Your --- request, I am sending you the observations I have been able to make of the new planets since then. All were made with the ring micrometer. Of these planets, only Vesta can be observed with our mural quadrant this year, and it still culminates in daylight.

\Needspace{9\baselineskip}
\begin{center}\emph{Observations of Juno.}\end{center}

\begin{center}\fontsize{10}{12}\selectfont\setlength{\tabcolsep}{3pt}\renewcommand{\arraystretch}{1.15}\begin{tabular}{lrrr}
1808 & \shortstack{Mean time\\in Göttingen} & \shortstack{Right\\ascension} & \shortstack{Southern\\declination}\\
June 20 & $11^{\text{u}}\,49^{\text{m}}\,0^{\text{s}}$ & $315^\circ\;29'\;34''$ & $2^\circ\;16'\;23''$\\
22 & $12\;0\;45$ & $315\;23\;1$ & $2\;14\;29$\\
July 6 & $12\;42\;23$ & $314\;1\;0$ & $2\;19\;7$\\
\end{tabular}\end{center}

\Needspace{9\baselineskip}\begin{center}\emph{Observations of Vesta.}\end{center}

\begin{center}\fontsize{9}{12}\selectfont\setlength{\tabcolsep}{3pt}\renewcommand{\arraystretch}{1.15}\begin{tabular}{lrrrl}
June 22 & $13^{\text{u}}\,46^{\text{m}}\,25^{\text{s}}$ & $353^\circ\;12'\;17''$ & $9^\circ\;12'\;16''$ & \\
July 1 & $14\;38\;57$ & $354\;28\;3$ & $9\;9\;39$ & \\
2 & $12\;50\;2$ & $355\;7\;7$ & & (by Prof. \textsc{Harding}).\\
13 & $13\;22\;37$ & $356\;35\;27$ & $9\;27\;40$ & \\
\end{tabular}\end{center}

\Needspace{7\baselineskip}\begin{center}\emph{Observations of Pallas.}\end{center}

\begin{center}\fontsize{9}{12}\selectfont\setlength{\tabcolsep}{3pt}\renewcommand{\arraystretch}{1.15}\begin{tabular}{lrrrl}
July 13 & $12^{\text{u}}\,3^{\text{m}}\,17^{\text{s}}$ & $300^\circ\;21'\;49''$ & $18^\circ\;41'\;44''$ north & \\
15 & $12\;1\;28$ & $299\;57\;3$ & $18\;32\;40$ & (by Prof. \textsc{Harding}).\\
\end{tabular}\end{center}

\par\indent Vesta has already reached the $7^{\text{th}}$ magnitude. Juno and Pallas, however, are very faint and can barely be judged to be of the $10^{\text{th}}$ magnitude. I have yet to see Ceres. The observations of Pallas will be of limited reliability, since they could only be made with a less powerful instrument, the planet now being too high for the Herschel telescope. --- For the elements I refer you to nos. 14, 40, and 107 of our \emph{Gelehrte Anzeigen} [pp. 292, 299, 300 of this volume]. I will prepare ephemerides for the coming year only after improving the elements again. For Vesta I see even less need for one, since once it disappears in March 1809 it will only be visible again in October or November 1809. --- Here are some more immersions of Jupiter’s satellites:

\begin{center}
\fontsize{10}{12}\selectfont\setlength{\tabcolsep}{3pt}\begin{tabular}{lccl}
$2^{\text{nd}}$ satellite June 22. & $12^{\text{u}}\,45^{\text{m}}\,9^{\text{s}}$ & & \textsc{Harding} $3\tfrac12$f. \textsc{Dollond}.\\
& $12\;45\;32$ & & \textsc{Gauss} $10$f. \textsc{Herschel}.\\
$1^{\text{st}}$ satellite July 13. & $13\;45\;36$ & & \textsc{Harding} $4$f. \textsc{Dollond}.\\
& $13\;45\;38$ & & \textsc{Thiarks} $3\tfrac12$f. \textsc{Dollond}.\\
& $13\;45\;38$ & & \textsc{Gauss} $10$f. \textsc{Herschel}.\\
\end{tabular}

Emersions.

\fontsize{10}{12}\selectfont\setlength{\tabcolsep}{3pt}\begin{tabular}{lccl}
$4^{\text{th}}$ satellite July 14. & $13\;35\;16$ & & \textsc{Harding} $10$f. \textsc{Herschel}.\\
& $13\;35\;56$ & & \textsc{Thiarks} $3\tfrac12$f. \textsc{Dollond}.\\
$2^{\text{nd}}$ satellite July 24. & $12\;32\;13$ & & \textsc{Thiarks} $3\tfrac14$f. \textsc{Dollond}.\\
& $12\;32\;27$ & & \textsc{Harding} $4$f. \textsc{Dollond}.\\
\end{tabular}
\end{center}

\Needspace{8\baselineskip}\par\indent Also 2 stellar occultations.

\smallskip
\noindent June 4.\quad $9^{\text{u}}\,3^{\text{m}}\,39^{\text{s}}0$\quad Immersion of $\iota$ Virginis. \textsc{Gauss} and \textsc{Harding} in the same second.

\smallskip
\noindent July 6.\quad $10\;39\;6.2$\quad Immersion of $\mu$ Sagittarii. \textsc{Gauss}.

\smallskip
{\rightskip=0pt plus 3em\par\indent The King has recently added to the instruments of our observatory by giving us a reflecting telescope by \textsc{Chevalier} and a timekeeper by \textsc{Berthoud}. --- The printing of my
\SourceLeaf{301}work on determining planetary orbits is still unfinished. Only 12 sheets are ready so far, and the whole will take about 30. We expect the $4^{\text{th}}$ sheet of Professor \textsc{Harding}’s charts by any post now. If only the public would give this valuable undertaking the support it needs! ---\par}
